\documentclass[11pt,a4paper]{ctexart}
\usepackage[margin=1in]{geometry}
\usepackage{booktabs}
\usepackage{longtable}
\usepackage{array}
\usepackage{hyperref}
\usepackage{enumitem}
\usepackage{xcolor}
\hypersetup{colorlinks=true,linkcolor=blue,urlcolor=blue}
\setlist[itemize]{nosep,leftmargin=1.5em}
\setlist[enumerate]{nosep,leftmargin=1.8em}
\title{Geant4-Agent Technical Route Manual\\Geant4-Agent ??????}
\author{}
\date{2026-03-08}
\begin{document}
\maketitle
\section*{English}
\subsection*{1. Objective}
\texttt{Geant4-Agent} is a controlled first-pass configuration agent for Geant4-oriented simulation setup. The immediate goal is not to generate a full production-ready Geant4 application, but to close the first-pass loop for five domains:
\begin{itemize}
\item geometry
\item materials
\item source
\item physics
\item output
\end{itemize}
The system accepts natural-language requests, normalizes them into structured proposals, validates them against deterministic rules, and commits them into an auditable configuration state. The language model does not write directly to the final config.
\subsection*{2. Core Technical Route}
The current route is built around five principles:
\begin{enumerate}
\item LLM-first for semantic normalization.
\item Single source of truth for committed configuration state.
\item Arbiter + validator gate as the only commit path.
\item Operator-graph geometry instead of per-scenario geometry templates.
\item Human approval for destructive overwrite.
\end{enumerate}
This is a control-plane architecture, not a pure prompting workflow.
\subsection*{3. Runtime Pipeline}
A single user turn is processed as follows:
\begin{enumerate}
\item The UI sends the current user message and the existing session state to the orchestrator.
\item The LLM normalization layer extracts intent, slot frame, semantic targets, and candidate updates.
\item BERT-based components provide conservative structure priors and fallback cues. They do not own the final decision.
\item Runtime semantic synthesis builds geometry candidates, graph programs, and grounded parameter proposals.
\item The arbiter merges user-grounded proposals, rejects unsafe conflicts, and enforces overwrite confirmation.
\item The validator gate checks schema completeness, semantic consistency, and feasibility constraints.
\item Only gated results are committed into the canonical config.
\item The dialogue layer renders a user-facing response, using action-level templates plus controlled LLM naturalization.
\end{enumerate}
\subsection*{4. Why the System Is Not LLM-Only}
A pure LLM flow failed for the expected reasons in this problem class:
\begin{itemize}
\item implicit overwrites across turns
\item geometry drift after partial clarification
\item internal field leakage into user-visible replies
\item hallucinated defaults that looked plausible but were not user-specified
\item inconsistent bindings between geometry, materials, and source settings
\end{itemize}
The current design therefore restricts the LLM to proposal generation. Commit authority remains in deterministic components.
\subsection*{5. Role of Each Major Layer}
\subsubsection*{5.1 LLM layer}
The LLM is responsible for semantic normalization, likely intent extraction, candidate field proposal, and limited naturalization of follow-up questions and summaries. It is not allowed to commit final configuration directly, override locked values without explicit user intent, or decide feasibility on its own.
\subsubsection*{5.2 BERT layer}
The BERT layer is intentionally downgraded from a central classifier to a supporting prior. It currently contributes structure prior, fallback semantic cues, and auxiliary confidence signals for routing.
\subsubsection*{5.3 Orchestrator and arbiter}
The orchestrator coordinates session state, pending overwrite flow, dialogue phase, and candidate lifecycle. The arbiter resolves candidate conflicts and ensures that only defensible updates survive to validation.
\subsubsection*{5.4 Validator gate}
The validator gate is the hard boundary before commit. It checks required-field completeness, path legality, semantic consistency, geometry feasibility rules, and cross-field coherence.
\subsubsection*{5.5 Dialogue layer}
The dialogue layer turns internal state into clarification questions, overwrite confirmation prompts, rejection acknowledgements, progress summaries, and final completion messages. Action templates are used here to keep outputs stable while preserving limited naturalization.
\subsection*{6. Geometry Route}
Geometry is the most specialized part of the stack. The route has two levels.
\subsubsection*{6.1 Single-solid families}
Current single-solid families include:
\begin{itemize}
\item \texttt{single\_box}
\item \texttt{single\_tubs}
\item \texttt{single\_sphere}
\item \texttt{single\_orb}
\item \texttt{single\_cons}
\item \texttt{single\_trd}
\item \texttt{single\_polycone}
\item \texttt{single\_cuttubs}
\item \texttt{single\_trap}
\item \texttt{single\_para}
\item \texttt{single\_torus}
\item \texttt{single\_ellipsoid}
\item \texttt{single\_elltube}
\item \texttt{single\_polyhedra}
\end{itemize}
\subsubsection*{6.2 Operator-graph families}
For composite layouts, the system uses graph-style assembly operators rather than per-scenario templates. Current graph families include:
\begin{itemize}
\item \texttt{ring}
\item \texttt{grid}
\item \texttt{nest}
\item \texttt{stack}
\item \texttt{shell}
\item \texttt{boolean}
\end{itemize}
This route is aligned with the original project requirement: finite assembly operators plus constraints, instead of scene-specific template classes.
\subsection*{7. Why Operator-Graph Matters}
A structure label alone does not scale well to real Geant4 requests. The operator-graph route generalizes better across scenarios with similar assembly semantics, and it supports partial completion and family-aware questioning without pretending the scene is already fully known.
\subsection*{8. Current Validation Strategy}
The system currently uses conservative validation rather than full Geant4 execution:
\begin{itemize}
\item parameter legality checks
\item containment checks
\item spacing checks for grid and ring
\item AABB-based sufficient conditions
\item structure-family completeness checks
\item overwrite and lock rules
\end{itemize}
This is intentional. The current agent is a first-pass validator, not yet the final Geant4 overlap judge.
\subsection*{9. Current Evidence}
Current controlled evaluation tracks:
\begin{itemize}
\item \texttt{v3\_24}: structured multi-turn live regression baseline
\item \texttt{v4\_24}: colloquial-material companion regression set
\end{itemize}
At the current checkpoint, both \texttt{v3\_24} and \texttt{v4\_24} have reached \texttt{24/24} in live regression. \texttt{v3\_24} remains the main baseline, while \texttt{v4\_24} is retained as the stricter colloquial-material companion set.
\subsection*{10. Known Remaining Gaps}
The main remaining gaps are coverage gaps, not architecture gaps:
\begin{itemize}
\item broader colloquial material grounding beyond the current casebank
\item more adversarial multi-turn stress cases
\item further operator-graph coverage under fuzzy language
\item larger public-facing regression bundles
\item tighter internal-leak scoring that measures user-visible leakage only
\end{itemize}
\subsection*{11. Recommended Next Steps}
\begin{enumerate}
\item Continue expanding the live casebank instead of rewriting the architecture.
\item Improve colloquial material grounding in the LLM-to-runtime handoff.
\item Keep BERT as a conservative prior unless a new gap clearly requires retraining.
\item Add more graph-heavy and adversarial demo cases before any open public trial.
\end{enumerate}
\newpage
\section*{??}
\subsection*{1. ??}
\texttt{Geant4-Agent} ????????? Geant4 ????????????????????????????? Geant4 ??????????????????????
\begin{itemize}
\item geometry
\item materials
\item source
\item physics
\item output
\end{itemize}
??????????????????????????????????????????????????????????????
\subsection*{2. ??????}
?????????????
\begin{enumerate}
\item ?? LLM-first ???????
\item ?????????????
\item Arbiter + Validator Gate ????????
\item ?????????????????????
\item ?????????????????
\end{enumerate}
????????? prompt ???????????????????
\subsection*{3. ??????}
??????????????
\begin{enumerate}
\item UI ???????????????? orchestrator?
\item LLM ?????? intent?slot frame?semantic targets ? candidate updates?
\item BERT ???????????????????????????
\item Runtime semantic ????????graph program ???????
\item Arbiter ???????????????? overwrite ?????
\item Validator Gate ? schema ??????????????????
\item ???? gate ??????? canonical config?
\item Dialogue ??????????????????????????????????
\end{enumerate}
\subsection*{4. ??????? LLM}
? LLM ??????????????????
\begin{itemize}
\item ?????????????
\item ??????????????
\item ?????????????
\item ?????????????????
\item geometry?materials?source ????????
\end{itemize}
?????????? LLM ??????????????????
\subsection*{5. ???????}
\subsubsection*{5.1 LLM ?}
LLM ??????????????????????????????????LLM ???????????????????????????????????????????
\subsubsection*{5.2 BERT ?}
BERT ??????????????????????? structure prior???????????????? routing ????????
\subsubsection*{5.3 Orchestrator ? Arbiter}
Orchestrator ?????????pending overwrite?phase ????????Arbiter ???????????????????????????
\subsubsection*{5.4 Validator Gate}
Validator Gate ???????????????????????????????????????????????????????
\subsubsection*{5.5 Dialogue ?}
Dialogue ????????????????????????????????????????????????????????????????????????
\subsection*{6. ??????}
?????????????????????
\subsubsection*{6.1 ?? solid family}
???????? family ???
\begin{itemize}
\item \texttt{single\_box}
\item \texttt{single\_tubs}
\item \texttt{single\_sphere}
\item \texttt{single\_orb}
\item \texttt{single\_cons}
\item \texttt{single\_trd}
\item \texttt{single\_polycone}
\item \texttt{single\_cuttubs}
\item \texttt{single\_trap}
\item \texttt{single\_para}
\item \texttt{single\_torus}
\item \texttt{single\_ellipsoid}
\item \texttt{single\_elltube}
\item \texttt{single\_polyhedra}
\end{itemize}
\subsubsection*{6.2 ????? family}
??????????? graph-style assembly operator????????????? graph family ???
\begin{itemize}
\item \texttt{ring}
\item \texttt{grid}
\item \texttt{nest}
\item \texttt{stack}
\item \texttt{shell}
\item \texttt{boolean}
\end{itemize}
???????????????????????????????????
\subsection*{7. ?????????}
????????????? Geant4 ??????????????????????????????????????? family-aware ????????????????
\subsection*{8. ??????}
?????????????????????? Geant4?
\begin{itemize}
\item ???????
\item containment ??
\item grid ? ring ?????
\item ?? AABB ?????
\item ????????
\item overwrite ? lock ??
\end{itemize}
????????????? first-pass validator?????? Geant4 overlap ???
\subsection*{9. ?????}
????????????
\begin{itemize}
\item \texttt{v3\_24}?????? live ????
\item \texttt{v4\_24}????????????
\end{itemize}
????????\texttt{v3\_24} ? \texttt{v4\_24} ???? live ????? \texttt{24/24}?\texttt{v3\_24} ???????\texttt{v4\_24} ????????????????
\subsection*{10. ??????}
???????????????????????
\begin{itemize}
\item ???????? grounding
\item ????????????
\item fuzzy ?????? operator-graph ??
\item ??????????
\item ????????????
\end{itemize}
\subsection*{11. ???????}
\begin{enumerate}
\item ??? live casebank???????????
\item ???? LLM ? runtime ??????????
\item ????????????? BERT ????????????????
\item ?????????????? graph-heavy ? adversarial ??????
\end{enumerate}
\end{document}
